\chapter{تحديد البنية: الرنين المغناطيسي النووي للكربون \texorpdfstring{\protect\babelsublr{$^{13}$C}}{13C} \omterm{def:b2:mass-spec-atomic:spectrum}{ومطيافية الكتلة} والأشعة تحت الحمراء معًا}\label{ch:b2:structure-determination}

يتلقى مختبر للعطور قارورة من سائل عديم اللون رائحته رائحة الياسمين. ولم يكتب أحد
ملصقًا. فيعطي \omterm{def:b2:mass-spec-atomic:spectrum}{طيف الكتلة} صيغته الجزيئية، وطيف الأشعة تحت الحمراء مجموعاته
الوظيفية، وطيف الكربون-13 هيكله وطيف البروتون طريقة اتصال القطع؛ وفي غضون
ساعة يصير للسائل اسم. ويضيف هذا الفصل \omterm{def:b2:structure-determination:c13}{الرنين المغناطيسي النووي للكربون-13} إلى الرنين المغناطيسي النووي للبروتون وإلى مطيافية الأشعة تحت الحمراء
في مجلد السنة الأولى وإلى \omterm{def:b2:mass-spec-atomic:spectrum}{مطيافية الكتلة} في \cref{ch:b2:mass-spec-atomic}، ويحوّل الأربعة
إلى طريقة واحدة.

\begin{recall}
مجلد السنة الأولى: الانزياح الكيميائي والحجب، والبروتونات المتكافئة، والتكامل، والاقتران سبين--سبين
وقاعدة $n + 1$، والأعداد الموجية للمجموعات الوظيفية في الأشعة تحت الحمراء ومنطقة البصمة، ودرجة
عدم التشبع، وحل بنية من ثلاثة أطياف. و\cref{ch:b2:mass-spec-atomic}: \omterm{def:b2:mass-spec-atomic:molecular-ion}{الأيون الجزيئي}،
\omterm{def:b2:mass-spec-atomic:isotope-pattern}{والأنماط النظائرية}، والتشظي.
\end{recall}

\section{الرنين المغناطيسي النووي للكربون-13}

ليس للكربون-12، النظير الرئيسي، سبين نووي ولا يعطي إشارة رنين مغناطيسي نووي. وللكربون-13 سبين
$\tfrac12$، مثل البروتون، لكنه لا يشكل إلا 1.1~\% من الكربون الطبيعي وإشارته أضعف،
نواةً بنواة، من إشارة البروتون؛ فيحتاج طيف الكربون إذن إلى عينة أكبر أو إلى مسوح أكثر.
% ledger: iso:C13

\begin{definition}[الرنين المغناطيسي النووي للكربون-13]\label{def:b2:structure-determination:c13}
\emph{الرنين المغناطيسي النووي للكربون-13}\index{رنين مغناطيسي نووي للكربون-13} يسجل رنينات نوى \ce{^{13}C} في عينة.
ويُجرى عادة مع \emph{نزع الاقتران العريض النطاق}\index{نزع الاقتران العريض النطاق}: إذ تُشعَّع
البروتونات على كامل مجال تردداتها أثناء القياس، مما يزيل كل
اقتران بين الكربون والبروتون، فيعطي كل نوع من الكربون خطًا واحدًا. و\emph{ذرات الكربون
المتكافئة}\index{ذرات كربون متكافئة}، التي يبادل بينها تناظر في الجزيء (أو دوران سريع)،
تعطي الخط نفسه.
\end{definition}

\begin{proposition}[لا اقتران بين الكربون والكربون]\label{prop:b2:structure-determination:no-cc-coupling}
لا تظهر الاقترانات بين ذرات الكربون المتجاورة في طيف \ce{^{13}C} بالوفرة
الطبيعية.
\end{proposition}

\begin{proof}
لا يُرى الاقتران بين ذرتي كربون إلا في الجزيئات التي تكون فيها الذرتان كلتاهما \ce{^{13}C}. ومن أجل موضعين
متجاورين معينين يكون الاحتمال $a_{13}^2 = 0.011^2 \approx 1.2 \times 10^{-4}$: أي نحو
جزيء واحد من ثمانية آلاف. وإشارة كل كربون تأتي كلها تقريبًا من جزيئات يكون فيها
جيرانه \ce{^{12}C}، غير المرئي للمطياف؛ والتوابع المقترنة أضعف مئة مرة
من ضجيج طيف عادي.
\end{proof}

\begin{proposition}[عدد الخطوط]\label{prop:b2:structure-determination:signals}
يُظهر طيف \ce{^{13}C} المنزوع الاقتران خطًا واحدًا لكل مجموعة من \omterm{def:b2:structure-determination:c13}{ذرات الكربون المتكافئة}، إلا إذا اتفق أن
كان لمجموعتين الانزياح نفسه.
\end{proposition}

\begin{proof}[تعليل]
\omterm{def:b2:structure-determination:c13}{ذرات الكربون المتكافئة} في محيطات متطابقة، فلها الانزياح نفسه؛ وغير المتكافئة
تختلف عمومًا. ومع إزالة الاقترانات مع البروتونات وغياب الاقترانات بين ذرات الكربون
(\cref{prop:b2:structure-determination:no-cc-coupling})، لا شيء يشطر الخط. ويحدث التراكب العرضي
حين يكون لكربونين مختلفين انزياحان أقرب من قدرة الميز، وهذا شائع
لذرات كربون حلقة البنزين قرب \qty{128}{ppm}.
\end{proof}

تنتشر الانزياحات على نحو \qty{220}{ppm}، أي عشرين ضعف مجال البروتونات، فيكون التراكب أندر منه
في أطياف البروتون، ويدل الانزياح وحده على نوع الكربون. ويبيّن المخطط أدناه الانزياحات
المقيسة لمركبات هذا الفصل.

\begin{omfigure}
\begin{tikzpicture}
\begin{axis}[width=0.86\linewidth, height=4.6cm, xmin=0, xmax=220, ymin=-0.6, ymax=4.6, x dir=reverse,
  xlabel={\babelsublr{$\delta$ (ppm)}}, label style={font=\small}, tick label style={font=\scriptsize},
  ytick={0,1,2,3,4}, yticklabels={\foreignlanguage{arabic}{كربون ألكيلي}, C--O, \foreignlanguage{arabic}{كربون عطري}, \foreignlanguage{arabic}{\babelsublr{C=O} إستر}, \foreignlanguage{arabic}{\babelsublr{C=O} كيتون}},
  xtick={0,20,40,60,80,100,120,140,160,180,200,220}, grid=major, grid style={black!10}]
\addplot[only marks, mark=|, mark size=5pt, thick, omDef] table[x=shift, y=cls] {figdata/out/bachelor-2-nmr-spectra-d.dat};
\end{axis}
\end{tikzpicture}

{\small انزياحات \ce{^{13}C} المقيسة للبيوتان-2-ون والبنتان-2-ون وبنزوات الإيثيل وإيثانوات البنزيل،
مرتبة بحسب نوع الكربون (CDCl$_3$؛ بيانات PubChem). كربونيلات الكيتونات قرب 209، وكربونيلات الإسترات قرب
167--171، والكربونات \omterm{def:b2:huckel:aromatic}{العطرية} بين 128 و137، والكربونات المرتبطة بأكسجين إستر قرب 61--66، 
والكربونات الألكيلية تحت 46.}
% figdata: bachelor-2/nmr-spectra.py part d
% ledger: c13:butanone, c13:pentan-2-one, c13:ethylbenzoate, c13:benzylacetate
\end{omfigure}

\begin{proposition}[الارتفاعات ليست أعدادًا]\label{prop:b2:structure-determination:intensities}
في طيف \ce{^{13}C} روتيني منزوع الاقتران لا تتناسب ارتفاعات الخطوط مع
أعداد ذرات الكربون التي تمثلها؛ والكربونات التي لا هيدروجين عليها تعطي خطوطًا ضعيفة.
\end{proposition}

\begin{proof}[تعليل]
تكرر الأطياف الروتينية القياس أسرع مما تسترخي أبطأ ذرات الكربون عائدة إلى التوازن، 
والكربونات التي لا بروتونات مرتبطة بها تسترخي ببطء، فتكون إشارتها مشبعة جزئيًا. ونزع الاقتران أيضًا
يعزز إشارات الكربونات الحاملة لبروتونات، بمقدار يتعلق بمحيطها.
ففي البيوتان-2-ون، خط الكربونيل أضعف الخطوط الأربعة مع أنه يمثل كربونًا واحدًا مثل
الأخرى. فالتكامل، المفيد جدًا للبروتونات، لا يُستعمل إذن في أطياف الكربون الروتينية.
\end{proof}

\section{تقنية DEPT}

\begin{definition}[تقنية DEPT]\label{def:b2:structure-determination:dept}
\emph{DEPT}\index{DEPT} (تعزيز غير مشوِّه بنقل الاستقطاب) مجموعة من تجارب الكربون-13
تعطي إشارات ذرات الكربون الحاملة لبروتونات إشارةً أو غيابًا بحسب
عدد هيدروجيناتها. و\emph{الكربون الرباعي}\index{كربون رباعي} لا يحمل هيدروجينًا
(مرتبط بأربع ذرات أخرى، أو كربون كربونيلي أو \omterm{def:b2:huckel:aromatic}{عطري} بلا H)؛ فلا يعطي إشارة DEPT.
\end{definition}

\begin{proposition}[قراءة DEPT]\label{prop:b2:structure-determination:dept}
في DEPT-135، تتجه إشارات كربونات \ce{CH} و\ce{CH3} إلى الأعلى، وكربونات \ce{CH2} إلى الأسفل والكربونات الرباعية
غائبة. وفي DEPT-90، لا تظهر إلا كربونات \ce{CH}.
\end{proposition}

\admitted

تنقل التجربة المغنطة من البروتونات إلى الكربون المرتبطة به، عبر
سلسلة من نبضات التردد الراديوي تختار زاويتها الأخيرة الاستجابة؛ وكيفية ذلك
تُشرح في مجلد السنة الثالثة. ومقارنة الطيف المنزوع الاقتران بطيفي DEPT-135 وDEPT-90 تصنّف كل
خط إلى C أو CH أو \ce{CH2} أو \ce{CH3}.

\begin{omfigure}
\pgfplotsset{dept/.style={width=\linewidth, height=2.9cm, ycomb, x dir=reverse, ymin=-1, ymax=1.15,
  ytick=\empty, axis y line=none, axis x line=middle, x axis line style={-}, tick label style={font=\tiny},
  title style={font=\scriptsize, at={(0.0,0.95)}, anchor=north west}, every axis plot/.append style={thick}}}
\begin{minipage}[t]{0.49\linewidth}\centering
{\small \foreignlanguage{arabic}{البيوتان-2-ون}}\par
\begin{tikzpicture}
\begin{axis}[dept, xmin=0, xmax=220, xtick={0,50,100,150,200}, title={\foreignlanguage{arabic}{منزوع الاقتران}}]
\addplot[omDef] table[x=shift, y=dec] {figdata/out/bachelor-2-nmr-spectra-a.dat};
\end{axis}
\end{tikzpicture}\par
\begin{tikzpicture}
\begin{axis}[dept, xmin=0, xmax=220, xtick={0,50,100,150,200}, title={\babelsublr{DEPT-135}}]
\addplot[omThm] table[x=shift, y=d135] {figdata/out/bachelor-2-nmr-spectra-a.dat};
\end{axis}
\end{tikzpicture}\par
\begin{tikzpicture}
\begin{axis}[dept, xmin=0, xmax=220, xtick={0,50,100,150,200}, title={\babelsublr{DEPT-90}}]
\addplot[omMeth] table[x=shift, y=d90] {figdata/out/bachelor-2-nmr-spectra-a.dat};
\end{axis}
\end{tikzpicture}
\end{minipage}\hfill
\begin{minipage}[t]{0.49\linewidth}\centering
{\small \foreignlanguage{arabic}{بنزوات الإيثيل}}\par
\begin{tikzpicture}
\begin{axis}[dept, xmin=0, xmax=220, xtick={0,50,100,150,200}, title={\foreignlanguage{arabic}{منزوع الاقتران}}]
\addplot[omDef] table[x=shift, y=dec] {figdata/out/bachelor-2-nmr-spectra-b.dat};
\end{axis}
\end{tikzpicture}\par
\begin{tikzpicture}
\begin{axis}[dept, xmin=0, xmax=220, xtick={0,50,100,150,200}, title={\babelsublr{DEPT-135}}]
\addplot[omThm] table[x=shift, y=d135] {figdata/out/bachelor-2-nmr-spectra-b.dat};
\end{axis}
\end{tikzpicture}\par
\begin{tikzpicture}
\begin{axis}[dept, xmin=0, xmax=220, xtick={0,50,100,150,200}, title={\babelsublr{DEPT-90}}]
\addplot[omMeth] table[x=shift, y=d90] {figdata/out/bachelor-2-nmr-spectra-b.dat};
\end{axis}
\end{tikzpicture}
\end{minipage}

{\small أطياف \ce{^{13}C} منزوعة الاقتران (انزياحات وارتفاعات مقيسة، بيانات PubChem) مع استجابات
\omterm{def:b2:structure-determination:dept}{DEPT} (الإشارات من \cref{prop:b2:structure-determination:dept}، والارتفاعات تخطيطية). البيوتان-2-ون: C=O
209.3 (غائب في \omterm{def:b2:structure-determination:dept}{DEPT})، و\ce{CH2} 36.9 (إلى الأسفل)، و\ce{CH3} 29.4 و7.9 (إلى الأعلى). بنزوات الإيثيل: C=O 166.5 
وكربون الحلقة الحامل للإستر، 130.6، يختفيان؛ وتبقى ثلاثة CH \omterm{def:b2:huckel:aromatic}{عطرية} (132.8، 129.6، 128.3) في
DEPT-90؛ و\ce{OCH2} 60.9 تتجه إلى الأسفل.}
% figdata: bachelor-2/nmr-spectra.py parts a, b
\end{omfigure}

\section{الجمع بين التقنيات}

كل تقنية تجيب عن سؤالها، وتُجمَّع الأجوبة بترتيب ثابت: الصيغة
أولًا، لأنها تحدّ كل ما عداها؛ ثم المجموعات الوظيفية والهيكل؛ ثم
الوصلات.

\begin{method}[حل مجهول]\label{met:b2:structure-determination:solve}
\begin{enumerate}
  \item الصيغة: \omterm{def:b2:mass-spec-atomic:molecular-ion}{الأيون الجزيئي}، و$M+1$ لعدد ذرات الكربون، \omterm{def:b2:mass-spec-atomic:isotope-pattern}{والأنماط النظائرية}، \omterm{def:b2:mass-spec-atomic:nitrogen-rule}{وقاعدة النيتروجين}؛ 
        والكتلة الدقيقة إن توافرت. ثم درجة عدم التشبع.
  \item الأشعة تحت الحمراء: \ce{C=O} (ومن أي نوع)، \ce{O-H}، \ce{N-H}، \ce{C#N}، و\ce{C-H} \omterm{def:b2:huckel:aromatic}{العطرية} أو الألكينية
        فوق \qty{3000}{cm^{-1}}.
  \item الكربون-13 \omterm{def:b2:structure-determination:dept}{وDEPT}: عدد أنواع الكربون (التناظر)، وأصنافها من الانزياحات، 
        وC وCH و\ce{CH2} و\ce{CH3} من \omterm{def:b2:structure-determination:dept}{DEPT}. وقارن بالصيغة.
  \item الرنين المغناطيسي النووي للبروتون: التكاملات، والانزياحات، والتعددات؛ والجيران من قاعدة $n + 1$.
  \item اجمع الشظايا؛ واكتب كل بنية مرشحة.
  \item تحقق من كل بنية مرشحة مقابل كل معطى، بما في ذلك شظايا \omterm{def:b2:mass-spec-atomic:spectrum}{طيف الكتلة}؛ واحتفظ بالتي
        تفسرها كلها.
\end{enumerate}
\end{method}

\begin{omfigure}
\begin{tikzpicture}[x=1cm, y=1cm, st/.style={draw=black!70, rounded corners=2pt, fill=white,
  font=\small, align=center, minimum height=0.9cm, inner sep=3pt, text width=2.3cm}]
\node[st, fill=omDef!10] (ms) at (0,0) {\babelsublr{MS}\\\foreignlanguage{arabic}{الصيغة، $M+1$}};
\node[st] (dbe) at (2.9,0) {درجة\\عدم التشبع};
\node[st, fill=omDef!10] (ir) at (5.8,0) {\babelsublr{IR}\\المجموعات الوظيفية};
\node[st, fill=omDef!10] (c13) at (8.7,0) {\ce{^{13}C}، \babelsublr{DEPT}\\الهيكل};
\node[st, fill=omDef!10] (h1) at (8.7,-1.8) {\ce{^1H}\\الوصلات};
\node[st] (cand) at (5.8,-1.8) {البنى\\المرشحة};
\node[st, fill=omThm!10] (chk) at (2.9,-1.8) {تحقق من كل\\معطى};
\node[st] (ans) at (0,-1.8) {البنية};
\draw[->, thick] (ms) -- (dbe); \draw[->, thick] (dbe) -- (ir); \draw[->, thick] (ir) -- (c13);
\draw[->, thick] (c13) -- (h1); \draw[->, thick] (h1) -- (cand); \draw[->, thick] (cand) -- (chk);
\draw[->, thick] (chk) -- (ans);
\draw[->, thick, densely dashed] (chk.south) -- ++(0,-0.5) -| node[pos=0.25, below, font=\footnotesize]
  {\foreignlanguage{arabic}{معطى غير مفسَّر: عُد إلى البنى المرشحة}} (cand.south);
\end{tikzpicture}
\par\vspace{4pt}

{\small ترتيب العمل في حالة مجهول (\cref{met:b2:structure-determination:solve}). والخطوة الأخيرة
هي الأكثر إغفالًا: فالبنية التي تترك قمة واحدة دون تفسير ليست الجواب بعد.}
\end{omfigure}

\section{حالات محلولة}

\begin{example}[كيتون، \babelsublr{C$_5$H$_{10}$O}]\label{ex:b2:structure-determination:ketone}
لمجهول $M = 86$، وحزمة في الأشعة تحت الحمراء قرب \qty{1715}{cm^{-1}}، وخطوط كربون عند 208.9 و45.7 و29.8 و17.4 
و13.7~ppm، وإشارات بروتون عند 2.41 (t، 2H)، و2.13 (s، 3H)، و1.61 (سداسية، 2H) و0.91 (t، 3H).
درجة عدم تشبع واحدة، وكيتون مشبع (لا بروتون ألدهيدي قرب 9.7). خمس ذرات كربون، وخمسة خطوط:
لا تناظر. والإشارة الأحادية لثلاثة بروتونات عند 2.13 هي \ce{CH3} مجاورة للكربونيل؛ والتسلسل ثلاثية،
سداسية، ثلاثية مجموعة بروبيل تلامس \ce{CH2} فيها عند 2.41 الكربونيل. إنه البنتان-2-ون،
\ce{CH3COCH2CH2CH3}. أما متماكبه المتناظر البنتان-3-ون فكان سيُظهر ثلاثة خطوط كربون ولا إشارة أحادية؛
ولطيف كتلته أيون ماكلافرتي عند 58 الذي يفتقر إليه البنتان-3-ون (\cref{ch:b2:mass-spec-atomic}).
% ledger: c13:pentan-2-one, nmr:pentan-2-one, ir:CO-ketone, ms:pentan-2-one
\end{example}

\begin{example}[إستر \omterm{def:b2:huckel:aromatic}{عطري}، \babelsublr{C$_9$H$_{10}$O$_2$}]\label{ex:b2:structure-determination:ester}
مجهول صيغته \ce{C9H10O2} (خمس درجات عدم تشبع) يُظهر حزمة كربونيل وسبعة
خطوط كربون: 166.5 (C)، و132.8 (CH)، و130.6 (C)، و129.6 (CH)، و128.3 (CH)، و60.9 (\ce{CH2}) و14.3
(\ce{CH3}). وطيف بروتونه: 8.05 (m، 2H)، و7.35--7.63 (m، 3H)، و4.37 (q، 2H)، و1.38 (t، 3H). 
وحلقة بنزين أحادية الاستبدال (أربعة خطوط كربون \omterm{def:b2:huckel:aromatic}{عطرية}، اثنان منها لذرتي كربون لكل منهما؛ وخمسة
بروتونات \omterm{def:b2:huckel:aromatic}{عطرية}) تفسر أربع درجات عدم تشبع، والكربونيل الخامسة. وزوج الرباعية--الثلاثية
مجموعة إيثيل؛ و\ce{CH2} فيها عند 4.37 (وعند 60.9 في طيف الكربون) مرتبطة بأكسجين. وكربونيل الإستر
عند 166.5 مرتبط بالحلقة: فالبروتونان عند 8.05، المجاوران له، مدفوعان نحو المجال المنخفض. إنه بنزوات
الإيثيل، \ce{C6H5COOCH2CH3}.
% ledger: c13:ethylbenzoate, nmr:ethylbenzoate
\end{example}

\section{أفخاخ}

\begin{itemize}
  \item \emph{البروتونات القابلة للتبادل} (\ce{O-H}، \ce{N-H}) تعطي إشارات عريضة عند انزياحات متغيرة، 
        وكثيرًا ما لا تقترن بجيرانها، وتختفي حين تُرج قطرة من \ce{D2O} مع
        العينة.
  \item \emph{الإشارات المتراكبة}: قد تتقاسم ذرتا كربون في حلقة بنزين، أو عدة مجموعات \ce{CH2} في سلسلة،
        خطًا واحدًا؛ فعُدّ الخطوط مقابل الصيغة قبل الاستنتاج بشأن التناظر.
  \item \emph{البروتونات اللامرآوية الموضع}: بروتونا \ce{CH2} المجاورة لمركز فراغي غير
        متكافئين، وقد يكون لهما انزياحان مختلفان ويقترن أحدهما بالآخر.
  \item \emph{الأنماط من الرتبة الثانية}: حين يكون لبروتونين مقترنين انزياحان أقرب من بضعة أضعاف
        ثابت اقترانهما، تميل المتعددات إحداها نحو الأخرى وتخفق قاعدة $n+1$؛ 
        والمطياف ذو المجال الأعلى يعيد الأنماط البسيطة.
\end{itemize}

\begin{safety}{\ghs[5mm]{GHS02}\,\ghs[5mm]{GHS07}\,\ghs[5mm]{GHS08}}
البيوتان-2-ون قابل للاشتعال. ومغناطيس مطياف الرنين المغناطيسي النووي يعمل دائمًا: فلا أدوات فولاذية، ولا أسطوانات
غاز ولا بطاقات مغناطيسية قربه، ولا دخول للأشخاص الحاملين لأجهزة تنظيم ضربات القلب.
% ledger: ghs:butanone, ghs:benzylacetate
\end{safety}

\section{تمارين}

\begin{exercise}[$\star$]\label{exo:b2:structure-determination:1}
كم خطًا يُظهر طيف \ce{^{13}C} المنزوع الاقتران لكل زايلين (1،2- و1،3- 
و1،4-ثنائي ميثيل بنزين)؟
\end{exercise}

\begin{exercise}[$\star$]\label{exo:b2:structure-determination:2}
تنبأ بطيفي DEPT-135 وDEPT-90 للبيوتان-2-ول.
\end{exercise}

\begin{exercise}[$\star$]\label{exo:b2:structure-determination:3}
إلى أي نوع من الكربون ينتمي على الأرجح خط عند 209، و170، و129، و64 و14~ppm؟
\end{exercise}

\begin{exercise}[$\star$]\label{exo:b2:structure-determination:4}
احسب درجة عدم التشبع للمركب \ce{C8H8O2} واقترح بنية تحوي حلقة بنزين.
\end{exercise}

\begin{exercise}[$\star\star$]\label{exo:b2:structure-determination:5}
سائل \ce{C4H8O} يُظهر حزمة أشعة تحت حمراء عند \qty{1715}{cm^{-1}} وخطوط كربون عند 209 و37 و29 و8~ppm
(DEPT-135: 37 إلى الأسفل، و29 و8 إلى الأعلى، و209 غائب). حدّد هويته.
\end{exercise}

\begin{exercise}[$\star\star$]\label{exo:b2:structure-determination:6}
كيف تميّز البنتان-2-ون من البنتان-3-ون بالرنين المغناطيسي النووي \ce{^{13}C} وحده؟ وبالرنين المغناطيسي النووي للبروتون؟ 
وبمطيافية الكتلة؟
\end{exercise}

\begin{exercise}[$\star\star$]\label{exo:b2:structure-determination:7}
لإستر $M = 88$ وإشارات بروتون عند 4.12 (q، 2H)، و2.04 (s، 3H) و1.26 (t، 3H) (معطيات
تمرين). حدّد هويته.
\end{exercise}

\begin{exercise}[$\star\star$]\label{exo:b2:structure-determination:8}
ميّز 1،2-ثنائي كلوروبنزين من 1،4-ثنائي كلوروبنزين بالرنين المغناطيسي النووي \ce{^{13}C}.
\end{exercise}

\begin{exercise}[$\star\star$]\label{exo:b2:structure-determination:9}
في 2-كلوروبيوتان، لماذا لا يكون بروتونا المجموعة \ce{CH2} متكافئين؟ وماذا يفعل ذلك
بطيف البروتون؟
\end{exercise}

\begin{exercise}[$\star\star\star$]\label{exo:b2:structure-determination:10}
مجهول \ce{C9H10O} يُظهر حزمة أشعة تحت حمراء عند \qty{1690}{cm^{-1}}، وخطوط كربون عند 200.8 (C)، و137.0 (C)،
و132.9 (CH)، و128.6 (CH)، و128.0 (CH)، و31.8 (\ce{CH2}) و8.3 (\ce{CH3})، وإشارات بروتون عند 7.95 (m،
2H)، و7.40--7.55 (m، 3H)، و2.98 (q، 2H) و1.22 (t، 3H) (معطيات تمرين). حدّد هويته، وفسّر
انخفاض العدد الموجي للكربونيل.
\end{exercise}

\begin{exercise}[$\star\star\star$]\label{exo:b2:structure-determination:11}
أربعة متماكبات \ce{C5H10O2}: حمض البنتانويك، وبيوتانوات الميثيل، وبروبانوات الإيثيل وإيثانوات البروبيل.
أعطِ لكل منها سمة واحدة من طيف الأشعة تحت الحمراء أو طيف البروتون تميزه من الثلاثة الأخرى.
\end{exercise}

\begin{exercise}[$\star\star\star$]\label{exo:b2:structure-determination:12}
ينسب تقرير الخط عند 166.5~ppm في بنزوات الإيثيل إلى كربون الحلقة الحامل للإستر 
والخط عند 130.6~ppm إلى كربون الكربونيل. بيّن، بمعطيات \omterm{def:b2:structure-determination:dept}{DEPT} ومخطط الانزياحات، أن هذا
النسب لا يمكن أن يكون صحيحًا.
\end{exercise}

\section{مسألة: أربعة أطياف للياسمين}

\begin{problem}[{مسألة نهاية الأسبوع --- الصيغة من الأيون الجزيئي، وحزمة الكربونيل في الأشعة تحت الحمراء،
والكربون-13 وDEPT، وطيف البروتون، وبنية يُتحقق منها مقابل الشظايا}]
\label{pb:b2:structure-determination:1}
سائل عديم اللون ذو رائحة زهرية كرائحة الياسمين يعطي المعطيات الآتية. \omterm{def:b2:mass-spec-atomic:spectrum}{طيف الكتلة} (قاعدة NIST
الشبكية): $m/z$ 150 (31.7)، 151 (3.1)، 108 (100)، 91 (71.7)، 90 (40.6)، 43 (37.6). الأشعة تحت الحمراء (معطيات
تمرين): حزمتان قويتان عند 1740 و\qty{1230}{cm^{-1}}، وحزم ضعيفة فوق \qty{3000}{cm^{-1}} بقليل، ولا
حزمة عريضة قرب \qty{3300}{cm^{-1}}. الكربون-13، CDCl$_3$ (بيانات PubChem): 170.70، 136.14، 128.56،
128.24، 66.24، 20.82~ppm. البروتون، CDCl$_3$، \qty{90}{MHz}: 7.33 (m، 5H)، 5.09 (s، 2H)، 2.06 (s، 3H).
% ledger: use:benzylacetate, ms:benzylacetate, c13:benzylacetate, nmr:benzylacetate, ir:CO-ester, iso:C12, iso:C13, mass:C12, mass:H1, mass:O16

\begin{center}
\begin{tikzpicture}
\pgfplotsset{dept/.style={width=0.8\linewidth, height=3.2cm, ycomb, x dir=reverse, ymin=0, ymax=1.15,
  xmin=0, xmax=200, ytick=\empty, axis y line=none, axis x line=bottom, x axis line style={-}, tick label style={font=\tiny},
  title style={font=\scriptsize, at={(0.02,0.82)}, anchor=west}, every axis plot/.append style={thick}}}
\begin{axis}[dept, title={\foreignlanguage{arabic}{منزوع الاقتران}}, xtick={0,20,40,60,80,100,120,140,160,180,200}]
\addplot[omDef] table[x=shift, y=dec] {figdata/out/bachelor-2-nmr-spectra-c.dat};
\end{axis}
\end{tikzpicture}
\end{center}
% figdata: bachelor-2/nmr-spectra.py part c

\medskip
\noindent\textbf{الجزء الأول --- الصيغة.}
\begin{enumerate}
  \item من النسبة 151/150، قدّر عدد ذرات الكربون.
  \item ماذا تقول \omterm{def:b2:mass-spec-atomic:nitrogen-rule}{قاعدة النيتروجين}؟
  \item مع تسع ذرات كربون، ماذا يبقى من الكتلة 150؟ اقترح صيغة فيها أكسجين.
  \item لماذا يكون \ce{C10H14O}، الذي كتلته الاسمية 150 أيضًا، أقل احتمالًا؟
  \item احسب درجة عدم التشبع.
  \item احسب \omterm{def:b2:mass-spec-atomic:exact-mass}{الكتلة الأحادية النظير} للصيغة المعتمدة.
\end{enumerate}

\medskip
\noindent\textbf{الجزء الثاني --- الأشعة تحت الحمراء.}
\begin{enumerate}[resume]
  \item انسب الحزمة عند \qty{1740}{cm^{-1}}. أي عائلة يوحي بها موضعها؟
  \item انسب الحزمة عند \qty{1230}{cm^{-1}}.
  \item علامَ تدل الحزم الضعيفة فوق \qty{3000}{cm^{-1}}؟
  \item ماذا يستبعد غياب حزمة عريضة قرب \qty{3300}{cm^{-1}}؟
  \item هل الكربونيل مترافق مع الحلقة؟ فسّر من عدده الموجي.
\end{enumerate}

\medskip
\noindent\textbf{الجزء الثالث --- الكربون-13 \omterm{def:b2:structure-determination:dept}{وDEPT}.}
\begin{enumerate}[resume]
  \item انسب كل خط إلى نوع من الكربون.
  \item أي خط يتجه إلى الأسفل في DEPT-135؟
  \item أي الخطوط تختفي في DEPT-135؟
  \item أي الخطوط تبقى في DEPT-90؟
  \item كم نوعًا من الكربون لحلقة بنزين أحادية الاستبدال؟ فسّر لماذا لا يظهر إلا خطان عطريان
        من نوع CH.
  \item الخط عند 128.24 هو الأطول. هل يمثل أكبر عدد من ذرات الكربون؟
\end{enumerate}

\medskip
\noindent\textbf{الجزء الرابع --- الرنين المغناطيسي النووي للبروتون والبنية.}
\begin{enumerate}[resume]
  \item انسب إشارات البروتون الثلاث.
  \item لماذا تكون الإشارتان عند 5.09 و2.06 أحاديتين؟
  \item لماذا تكون \ce{CH2} عند 5.09 بعيدة جدًا نحو المجال المنخفض؟
  \item اجمع البنية وسمِّها.
  \item لمتماكبها 2-فينيل إيثانوات الميثيل، \ce{C6H5CH2COOCH3}، الصيغة نفسها. أي معطى
        يستبعده؟
  \item فسّر الشظايا عند 108 و91 و43.
  \item كم خطًا كان سيُظهر طيف \ce{^{13}C} منزوع الاقتران تام الميز للمركب، 
        وكم منها يتجه إلى الأسفل في DEPT-135؟
\end{enumerate}
\end{problem}
